% Generated by scripts/08_numbers.R - do not edit by hand.
\newcommand{\StationDataText}{Se usaron 63 estaciones con observaciones horarias entre el 1 de enero y el 31 de marzo de 2019 dentro del rectángulo \ang{34}--\ang{38.5}~S, \ang{70.3}--\ang{73.6}~O (figura~\ref{fig:area}): 5 de la Dirección Meteorológica de Chile (DMC), 21 de la red del Instituto de Investigaciones Agropecuarias (INIA) y 35 de la Fundación para el Desarrollo Frutícola (FDF), todas obtenidas del portal público SACLIM de la DMC, más 2 estaciones de aeropuerto de la base NOAA ISD no repetidas en la DMC. Todas registran temperatura y humedad del aire (punto de rocío o humedad relativa, de la que se obtiene la presión de vapor con FAO-56); solo 7 miden viento, a \SI{10}{\metre}. El portal publica la radiación global como máximo horario de lecturas instantáneas, no como media, por lo que no se usa como forzamiento (D-04). Los valores a HH:00~UTC se tratan como puntuales. Las alturas de los sensores de temperatura se asumieron en \SI{2}{\metre}.}
% Affiliations as given by the author in his CEDES 2027 extended abstract (desierto florido).
\newcommand{\AuthorAffiliation}{$^{1}$Instituto de Ciencias y Tecnología Ambiental (ICTA Ltda.), Chile; $^{2}$EARTH Lab (Environmental Analysis, Remote Sensing \& Terrestrial Hydrology Laboratory), Universidad San Sebastián, Chile}
